\SGANextSectionTarget{sga3:xxiii:section:6}\SGASection{6}{Chevalley Systems}

The explicit calculations of \SGARefLink{sga3:xxiii:section:3}{Section 3} have important numerical
consequences.

\SGARefTarget{sga3:XXIII:6.1}
\SGAStatement{Definition 6.1}{Let \(S\) be a scheme and
\((G,T,M,R)\) a split \(S\)-group. A \emph{Chevalley system} of
\(G\) is a family
\[
  (X_\alpha)_{\alpha\in R},
  \qquad
  X_\alpha\in\Gamma(S,\mathfrak g^\alpha)^\times,
\]
satisfying:
\begin{description}
  \item[(SC)] for every pair \(\alpha,\beta\in R\),
  \[
    \operatorname{Ad}(w_\alpha(X_\alpha))X_\beta
      =\pm X_{s_\alpha(\beta)}.
  \]
\end{description}}

Recall that
\[
  w_\alpha(X_\alpha)
    =\exp(X_\alpha)\exp(-X_\alpha^{-1})\exp(X_\alpha).
\]
Condition (SC) implies, in particular,
\[
  X_{-\alpha}=\pm X_\alpha^{-1}
\]
by the identity
\[
  \operatorname{Ad}(w_\alpha(X_\alpha))X_\alpha=-X_\alpha^{-1}.
\]

\SGARefTarget{sga3:XXIII:6.2}
\SGAStatement{Proposition 6.2}{Every split group admits a Chevalley
system. More precisely, if
\[
  \bigl(\Delta,(X'_\alpha)_{\alpha\in\Delta}\bigr)
\]
is a pinning of \((G,T,M,R)\), there is a Chevalley system
\((X_\alpha)_{\alpha\in R}\) with \(X_\alpha=X'_\alpha\) for
\(\alpha\in\Delta\).}

First, it suffices to verify (SC) for \(\alpha\in\Delta\). Indeed,
for every \(\alpha\in R\), choose a sequence
\(\alpha_i\in\Delta\) such that
\[
  \alpha=s_{\alpha_1}\cdots s_{\alpha_n}(\alpha_{n+1}).
\]
Repeated use of (SC) gives
\[
 X_\alpha
  =\pm\operatorname{Ad}\!\left(
    w_{\alpha_1}(X_{\alpha_1})\cdots
    w_{\alpha_n}(X_{\alpha_n})\right)X_{\alpha_{n+1}}.
\]
By \SGARefLink{sga3:xxiii:lemma:3-1-1:item:iii}{Lemma 3.1.1(iii)},
\[
\begin{aligned}
 &w_{\alpha_1}(X_{\alpha_1})\cdots w_{\alpha_n}(X_{\alpha_n})
  w_{\alpha_{n+1}}(X_{\alpha_{n+1}})
  w_{\alpha_n}(X_{\alpha_n})^{-1}\cdots
  w_{\alpha_1}(X_{\alpha_1})^{-1}\\
 &\hspace{42mm}=\alpha^*(\pm1)w_\alpha(X_\alpha).
\end{aligned}
\]
Now
\[
 w_{\alpha_i}(X_{\alpha_i})^{-1}
  =\alpha_i^*(-1)w_{\alpha_i}(X_{\alpha_i}),
\]
and
\[
 \beta(\gamma^*(-1))=(-1)^{(\gamma^*,\beta)}=\pm1.
\]
It follows that (SC) for every pair \((\alpha_i,\gamma)\) implies
(SC) for every \((\alpha,\beta)\).

For each \(\alpha\in R\), choose such a sequence and define
\[
 X_\alpha=
 \operatorname{Ad}\!\left(
  w_{\alpha_1}(X_{\alpha_1})\cdots
  w_{\alpha_n}(X_{\alpha_n})\right)X'_{\alpha_{n+1}}.
\]
It remains to prove the following lemma.

\SGARefTarget{sga3:XXIII:6.3}
\SGAStatement{Lemma 6.3}{Let
\[
 (G,T,M,R,\Delta,(X_\alpha)_{\alpha\in\Delta})
\]
be a pinned \(S\)-group. If \(\alpha_i\), \(0\le i\le n+1\), are
simple roots such that
\[
  s_{\alpha_1}\cdots s_{\alpha_n}(\alpha_{n+1})=\alpha_0,
\]
then
\[
  \operatorname{Ad}(w_{\alpha_1}\cdots w_{\alpha_n})
    X_{\alpha_{n+1}}=\pm X_{\alpha_0}.
\]}

Arguing as in \SGARefLink{sga3:XXIII:2.3.4}{2.3.4}, using Tits's lemma, reduces the proof to:
\begin{enumerate}
  \item[\SGAInlineRefTarget{sga3:XXIII:6.3:item:a}(a)] \(G\) has semisimple rank at most \(2\); or
  \item[\SGAInlineRefTarget{sga3:XXIII:6.3:item:b}(b)] \(w_{\alpha_1}\cdots w_{\alpha_n}\) is a section of
    \(T\).
\end{enumerate}
Case \SGARefLink{sga3:XXIII:6.3:item:a}{(a)} is exactly the calculation in \SGARefLink{sga3:XXIII:3.1.2:item:i}{3.1.2(i)}, \SGARefLink{sga3:xxiii:proposition:3-2-1:item:i}{3.2.1(i)},
\SGARefLink{sga3:xxiii:proposition:3-3-1:item:i}{3.3.1(i)}, or \SGARefLink{sga3:xxiii:proposition:3-4-1:item:i}{3.4.1(i)}.

In case \SGARefLink{sga3:XXIII:6.3:item:b}{(b)}, it suffices to show that if
\[
  s_{\alpha_1}\cdots s_{\alpha_n}=1,
\]
then \(t=w_{\alpha_1}\cdots w_{\alpha_n}\) satisfies
\(\alpha(t)=\pm1\) for every root \(\alpha\). By the structure of the
Weyl group, the word in the free group generated by the
\(s_\alpha\) lies in the normal subgroup generated by
\[
  (s_\alpha s_\beta)^{n_{\alpha\beta}}.
\]
Thus one reduces to a word
\[
 s_{\alpha_1}\cdots s_{\alpha_i}
 (s_{\alpha_{i+1}}s_{\alpha_{i+2}})
 ^{n_{\alpha_{i+1}\alpha_{i+2}}}
 s_{\alpha_i}\cdots s_{\alpha_1},
\]
for which
\[
  t=s_{\alpha_1}\cdots s_{\alpha_i}
      (t_{\alpha_{i+1}\alpha_{i+2}}).
\]
The assertion is now immediate from the values of
\(t_{\alpha_1\alpha_2}\) computed in \SGARefLink{sga3:XXIII:3.1.2:item:i}{3.1.2(i)}, \SGARefLink{sga3:xxiii:proposition:3-2-1:item:i}{3.2.1(i)},
\SGARefLink{sga3:xxiii:proposition:3-3-1:item:i}{3.3.1(i)}, and \SGARefLink{sga3:xxiii:proposition:3-4-1:item:i}{3.4.1(i)}.

\SGARefTarget{sga3:XXIII:6.4}
\SGAStatement{Proposition 6.4}{Let \((G,T,M,R)\) be a split
\(S\)-group with Chevalley system \((X_\alpha)_{\alpha\in R}\).
Let \(\alpha,\beta\) be nonproportional roots, with
\[
  \operatorname{length}(\alpha)\le\operatorname{length}(\beta),
  \qquad
  |(\beta^*,\alpha)|\le|(\alpha^*,\beta)|.
\]
Let \(q\) and \(p-1\) be the nonnegative integers for which the roots
of the form \(\beta+k\alpha\) are
\[
  \{\beta-(p-1)\alpha,\ldots,\beta,\ldots,\beta+q\alpha\}.
\]
Thus \(-(\alpha^*,\beta)=q-p+1\). Writing
\(p_\gamma(x)=\exp(xX_\gamma)\), the commutation relation between
\(U_\alpha\) and \(U_\beta\) is given by:
\[
\begin{array}{c|l}
(p,q)&p_\beta(y)p_\alpha(x)p_\beta(-y)p_\alpha(-x)\\ \hline
(-,0)&e\\
(1,1)&p_{\alpha+\beta}(\pm xy)\\
(1,2)&p_{\alpha+\beta}(\pm xy)p_{2\alpha+\beta}(\pm x^2y)\\
(1,3)&p_{\alpha+\beta}(\pm xy)p_{2\alpha+\beta}(\pm x^2y)
       p_{3\alpha+\beta}(\pm x^3y)p_{3\alpha+2\beta}(\pm x^3y^2)\\
(2,1)&p_{\alpha+\beta}(\pm2xy)\\
(2,2)&p_{\alpha+\beta}(\pm2xy)p_{2\alpha+\beta}(\pm3x^2y)
       p_{\alpha+2\beta}(\pm3xy^2)\\
(3,1)&p_{\alpha+\beta}(\pm3xy)
\end{array}
\]
These are all possibilities, since the root string has length at
most \(3\).}

\SGAProof
By \SGARefLink{sga3:XXI:3.5.4}{Exposé XXI, 3.5.4}, there is a system of simple roots \(\Delta\)
such that \(\alpha\in\Delta\), and there are
\(\alpha'\in\Delta\), \(a,b\in\mathbf Q_+\), with
\[
  \beta=a\alpha+b\alpha'.
\]
The desired relation lies in \(U_{\alpha\alpha'}\), so it follows
from the rank-two calculations of Section 3 and condition (SC).

\SGARefTarget{sga3:XXIII:6.5}
\SGAStatement{Corollary 6.5 (Chevalley's rule)}{Let
\((X_\alpha)_{\alpha\in R}\) be a Chevalley system of the split
\(S\)-group \(G\). If \(\alpha,\beta,\alpha+\beta\in R\), then
\[
  [X_\alpha,X_\beta]=\pm pX_{\alpha+\beta},
\]
where \(p\) is the least positive integer for which
\(\beta-p\alpha\) is not a root.}

The assertion is symmetric in \(\alpha,\beta\); after arranging
\(\operatorname{length}(\alpha)\le\operatorname{length}(\beta)\),
it follows from \SGARefLink{sga3:XXIII:6.4}{Proposition 6.4}.

\SGARefTarget{sga3:XXIII:6.6}
\SGAStatement{Corollary 6.6}{Let \(S\) be a scheme such that
\(6\cdot1_S\ne0\), and let \((G,T,M,R)\) be a split \(S\)-group.
If \(R'\subset R\) is such that
\begin{equation}
  \mathfrak g_{R'}
   =\mathfrak t\oplus\bigoplus_{\alpha\in R'}\mathfrak g^\alpha
  \tag{*}\label{sga3:XXIII:6.6:star}
\end{equation}
is a Lie subalgebra of \(\mathfrak g\), then \(R'\) is a closed
subset of \(R\).}

Choose a Chevalley system and let \(\alpha,\beta\in R'\) with
\(\alpha+\beta\in R\). By \SGARefLink{sga3:XXIII:6.5}{Corollary 6.5} and \SGARefLink{sga3:XXIII:6.4}{Proposition 6.4},
\[
  [X_\alpha,X_\beta]=\pm pX_{\alpha+\beta},
  \qquad p\in\{1,2,3\}.
\]
Since neither \(2\) nor \(3\) is zero on \(S\),
\eqref{sga3:XXIII:6.6:star} implies
\(\mathfrak g^{\alpha+\beta}\subset\mathfrak g_{R'}\), and hence
\(\alpha+\beta\in R'\).
\NDE{It suffices that \(2\) and \(3\) are nonzero on \(S\); for
example the result applies to \(S=\operatorname{Spec}(\mathbf Z/6\mathbf Z)\).
In the excluded characteristics, short-root subsystems furnish
standard counterexamples.}

\SGARefTarget{sga3:XXIII:6.7}
\SGAParagraph{6.7}
The exact signs in this section can be determined by studying
\(\operatorname{Norm}_G(T)\), and more precisely the extended Weyl
group
\[
  W^*=N(\mathbf Z),\qquad
  N=\operatorname{Norm}_{\operatorname{\acute{E}p}_{\mathbf Z}(R)}
         (T_{\mathbf Z}(R)).
\]
It is an extension of \(W(R)\) by an abelian group of type
\((2,2,\ldots,2)\), which accounts for the signs.

\SGARefTarget{sga3:XXIII:6.7.1}
\SGAStatement{Remark 6.7.1}{By 3.1.2(i) and 3.\(n\).1,
\(n=2,3,4\), the elements \(w_\alpha,w_\beta\in N(\mathbf Z)\)
satisfy the braid relations
\[
  w_\alpha w_\beta\cdots=w_\beta w_\alpha\cdots
\]
with \(n_{\alpha\beta}\) factors on each side.}

See also Tits \cite{Ti66}, Bourbaki \cite{BLie}, IX.4, Ex.~12, and
Springer \cite[9.3.2]{Sp98}.

\begin{thebibliography}{99}
\bibitem[BLie]{BLie}
N.~Bourbaki,
\emph{Groupes et algèbres de Lie}, Chap.~IX,
Masson, 1982.

\bibitem[Ja87]{Ja87}
J.~C. Jantzen,
\emph{Representations of Algebraic Groups},
Academic Press, 1987; 2nd ed., American Mathematical Society, 2003.

\bibitem[Sp98]{Sp98}
T.~A. Springer,
\emph{Linear Algebraic Groups}, 2nd ed.,
Birkhäuser, 1998.

\bibitem[St62]{St62}
R.~Steinberg,
\emph{Générateurs, relations et revêtements de groupes algébriques},
Colloque sur la théorie des groupes algébriques (Bruxelles, 1962),
Université de Louvain and Gauthier--Villars, 1962;
pp.~133--147 in \emph{R. Steinberg Collected Papers},
American Mathematical Society, 1997.

\bibitem[St67]{St67}
R.~Steinberg,
\emph{Lectures on Chevalley Groups},
Yale University, 1967.

\bibitem[Ta83]{Ta83}
M.~Takeuchi,
\emph{A hyperalgebraic proof of the isomorphism and isogeny
theorems for reductive groups},
J. Algebra \textbf{85} (1983), 179--196.

\bibitem[Ti66]{Ti66}
J.~Tits,
\emph{Normalisateurs de tores I. Groupes de Coxeter étendus},
J. Algebra \textbf{4} (1966), 96--116.
\end{thebibliography}
